\documentclass[12pt]{article}

% Circuito com 6 resistencias
% Gambelas, Ter 10 Dez 2019 14:43:38 WET 
% Written by Tordar

\usepackage{graphicx}
\usepackage{pst-circ}

\pagestyle{empty}

\begin{document}

\begin{pspicture}(12,5)

\pnode(0,2){P1}
\pnode(5,2){P2}
\pnode(8,2){P3}
\pnode(8,-2){P4}
\pnode(5,-2){P5}
\pnode(0,-2){P6}

\resistor[labeloffset=1,dipolestyle=zigzag](P1)(P2){\scalebox{1.5}{$R_1$}}
\resistor[labeloffset=1,dipolestyle=zigzag](P2)(P3){\scalebox{1.5}{$R_2$}}
\resistor[labeloffset=1,dipolestyle=zigzag](P3)(P4){\scalebox{1.5}{$R_3$}}
\resistor[labeloffset=1,dipolestyle=zigzag](P4)(P5){\scalebox{1.5}{$R_4$}}
\resistor[labeloffset=1,dipolestyle=zigzag](P5)(P6){\scalebox{1.5}{$R_5$}}
\resistor[labeloffset=1,dipolestyle=zigzag](P2)(P5){\scalebox{1.5}{$R_6$}}

\pscircle[fillstyle=solid,fillcolor=white](0,2){0.15}
\pscircle[fillstyle=solid,fillcolor=white](0,-2){0.15}

\end{pspicture}

\end{document}
